\medskip 
{\color{gray}\hrule}
\section{General Decomposition in Degree $\geq 3$}
\medskip 
{\color{gray}\hrule}
\medskip
Let $K$ be a number field of degree $n \geq 3$, and let $p$ be a prime number. Inside $\mathcal{O}_K$, the ideal generated by $p$ no longer has to remain prime. Instead, it factors as
\[
(p) = \mathfrak{p}_1^{e_1} \cdots \mathfrak{p}_g^{e_g},
\]
where the $\mathfrak{p}_i$ are distinct prime ideals.\\
Each of these prime ideals comes with two numerical measures. First, the exponent $e_i$, is called the $ramification$ $index$. It records how many times $\mathfrak{p}_i$ appears in the factorization. Second, the $inertia$ $degree$
\[
f_i = [\mathcal{O}_K/\mathfrak{p}_i : \mathbb{F}_p]
\].\\
These numbers are not independent as they satisfy the relation
\[
\sum_{i=1}^g e_i f_i = [K : \mathbb{Q}] = n.
\]
This identity can be seen as a kind of balance: each prime above $p$ contributes an amount $e_i f_i$ to the total degree, and together they add up to the whole degree of the extension. In particular, this strongly limits how the prime can factor.\\
Another useful perspective comes from norms. For each prime ideal $\mathfrak{p}_i$, then
\[
N(\mathfrak{p}_i) = p^{f_i}.
\]
Using multiplicativity of the norm, this gives
\[
p^n = N((p)) = \prod_{i=1}^g p^{e_i f_i},
\]
which recovers the same identity.\\\\
In degree $2$, the possibilities are limited; 
\[
2 = 1+1,\quad 2 = 2
\]
For extensions of degree at least $3$, we start to see different types of behavior. 
For example, in degree $4$, all possible factorizations correspond to the partitions of $4$:
\[
4 = 1+1+1+1, \quad 4 = 2+1+1, \quad 4 = 2+2, \quad 4 = 3+1, \quad 4 = 4,
\]
here ramification allows repeated terms.\\
In higher degrees, there are more such partitions, and each one corresponds to a different way a prime can factor in the extension.

